% ============================================================================
% Section 1: Introduction
% ============================================================================
\section{Introduction}
\label{sec:introduction}

\subsection{Historical Context}

Srinivasa Ramanujan (1887--1920) left behind three mathematical
notebooks and a collection of loose papers---the so-called
\emph{Lost Notebook}---containing thousands of formulas, most
stated without proof~\cite{ramanujan_lost,andrews_berndt_I}.
A significant fraction of these results involve the Dedekind
eta function
\begin{equation}
  \label{eq:dedekind_eta}
  \eta(\tau) = q^{1/24} \prod_{n=1}^{\infty} (1 - q^n),
  \qquad q = e^{2\pi i \tau},\quad \mathrm{Im}(\tau) > 0,
\end{equation}
and products of the form
$\prod_{d \mid N} \eta(d\tau)^{r_d}$, now called
\emph{eta-quotients}~\cite{ono2004,martin_1996}.
Despite over a century of scholarship, a systematic
computational recovery of all such identities latent
in the original manuscripts has never been attempted.

\subsection{Contribution}

This paper presents the \emph{Ramanujan Autonomous
Neuro-symbolic Architecture} (\RAMA{}), a five-stage
discovery pipeline that:
\begin{enumerate}[label=(\roman*)]
  \item extracts mathematical content from digitized
        manuscript page images via multimodal vision models;
  \item searches the combinatorial space of eta-quotient
        exponent vectors $(r_1, r_2, \ldots, r_{d_{\max}})$
        using a population-based genetic algorithm with
        tournament selection;
  \item evaluates candidate fitness against target
        $q$-series coefficients using a composite energy
        functional $E = \alpha C + \beta I + \gamma D$;
  \item auto-formalizes each candidate as a Lean~4 theorem
        and certifies it via the Lean kernel (Mathlib);
  \item maps verified quotients to physical interpretations
        via saddle-point thermodynamics.
\end{enumerate}

\subsection{Summary of Results}

As of this preliminary report, processing 481 of 698
manuscript pages from Notebooks~1, 2, and~3:
\begin{itemize}
  \item \textbf{477 candidate eta-quotients} have been
        formally verified by the Lean~4 kernel
        (\texttt{sorry}-free, zero axiomatic gaps).
  \item \textbf{378} of these correspond to thermodynamically
        \emph{stable} vacuum states ($\ceff > 0$).
  \item The best-fit candidate achieves a RAMA energy of
        $E = 1.1272$ with a fit error of $I = 0.0705$
        (7.05\% residual against target coefficients).
\end{itemize}

\subsection{Epistemic Disclaimer}

We emphasize that the results reported here are
\emph{computational discoveries}: conjectures identified
by a search algorithm and verified at the level of
algebraic type-checking by Lean~4.  The formal verification
confirms that the stated Lean code compiles without
\texttt{sorry} or axiomatic escape hatches, but it does
\emph{not} constitute a proof that the discovered
eta-quotients match the original intent of Ramanujan's
manuscript entries.  Cross-referencing with established
results in the literature~\cite{andrews_berndt_I,andrews_berndt_II,ono2004}
is an ongoing effort.

\subsection{Paper Organization}

Section~\ref{sec:math_prelim} reviews the mathematical
background on eta-quotients, modular forms, and mock theta
functions.
Section~\ref{sec:pipeline} describes the \RAMA{} pipeline
architecture.
Sections~\ref{sec:candidate_alpha}--\ref{sec:candidate_gamma}
present three representative discoveries with full proofs.
Section~\ref{sec:statistics} provides a statistical analysis
of the full corpus.
Section~\ref{sec:physics} discusses physical interpretations
and their limitations.
Section~\ref{sec:limitations} catalogs known limitations and
falsifiability conditions.
Section~\ref{sec:conclusion} concludes.
Appendix~\ref{app:lean4} contains the complete Lean~4 code;
Appendix~\ref{app:python} contains Python reproducibility
scripts.
